\uuid{1Yn3}
\titre{Points stationnaires}
\auteur{hanani ; mzari}
\organisation{exo7}
\licence{CC BY-NC-SA 4.0}
\datecreate{2020-01-20}
\niveau{L1}
\module{Géométrie}
\chapitre{Courbes planes}
\sousChapitre{Courbes paramétrées}
\difficulte{3}
\theme{Courbes paramétrées, points stationnaires}
\exo7qcmid{lille-2.courbes.007}

\contenu{
\begin{question}
\qclassification{244.01}
\qtags{motcle=points stationnaires, tangentes}
La trajectoire $\Gamma$ d'une particule en mouvement est donnée par les équations
$$x=2+\cos t\quad \mbox{et} \quad y=\frac{t^2}{2}+\sin t.$$
\begin{answers}
\good{$\Gamma$ n'admet pas de point stationnaire.}
\good{Le point de paramètre $t=\pi/2$ est un point d'inflexion.}
\bad{La tangente à $\Gamma$ au point de paramètre $t=\pi/2$ est verticale.}
\bad{$\Gamma$ est symétrique par rapport à l'axe des $x$.}
\end{answers}
\begin{explanations}
Notons $f(t)=(x(t),y(t))$. Le système $f'(t)=(0,0)$ n'admet pas de solution, donc $\Gamma$ n'admet pas de point stationnaire. $\displaystyle \lim _{t\to \pm\infty}x(t)$ n'existe pas. Un DL de f en $\pi/2$ montre que le point de paramètre $t=\pi/2$ est un point d'inflexion et la tangente en ce point est dirigée par le vecteur $f'(\pi/2)=(-1,\pi/2)$. Enfin $y$ n'est ni paire ni impaire, donc $\Gamma$ n'est pas symétrique par rapport à l'axe des $x$.
\end{explanations}
\end{question}
}
